\documentclass[11pt]{amsart}

\usepackage[T1]{fontenc}
\usepackage{lmodern}
\usepackage{microtype}
\usepackage{amsmath,amssymb,amsthm}
\usepackage[hidelinks]{hyperref}

\hypersetup{
  pdftitle={The n=20 Case of the Induced Planar P6-Turan Conjecture}
}

\newtheorem{theorem}{Theorem}
\newtheorem{proposition}[theorem]{Proposition}
\newtheorem{lemma}[theorem]{Lemma}

\newcommand{\cP}{\mathcal P}
\newcommand{\cL}{\mathcal L}
\newcommand{\exP}{\operatorname{ex}_{\cP}}
\newcommand{\ind}{\mathrm{ind}}

\title[The induced planar $P_6$-Tur\'an conjecture at $n=20$]
{The $n=20$ Case of the Induced Planar $P_6$-Tur\'an Conjecture}

\date{}

\subjclass[2020]{05C35, 05C10}
\keywords{induced planar Tur\'an number, induced path, planar triangulation,
exhaustive computation}

\begin{document}

\begin{abstract}
Gy\H{o}ri and Hama Karim conjectured an exact formula for the induced planar
Tur\'an number of the six-vertex path. We verify its first case below
Euler's planar bound:
\[
  \exP(20,P_6^{\ind})=53.
\]
The lower bound is their construction. The upper bound follows from a
computer-assisted census: the unique induced-$P_6$-free triangulation on
$19$ vertices has no induced-$P_6$-free facial extension, and no
$20$-vertex triangulation of minimum degree at least four is
induced-$P_6$-free.
\end{abstract}

\maketitle

\section{The exact value}

All graphs are finite and simple. A triangulation is a plane graph in which
every face is bounded by a triangle. For a graph $F$, let
$\exP(n,F^{\ind})$ denote the maximum number of edges in an $n$-vertex
planar graph containing no induced copy of $F$.

Gy\H{o}ri and Hama Karim \cite[Conjecture~5.1]{GyoriKarim} conjectured
that, for every $n\geq 3$,
\begin{equation}\label{eq:conjecture}
\exP(n,P_6^{\ind})=
\begin{cases}
3n-6, & n\leq 19,\\[1mm]
\left\lfloor \dfrac{50}{17}(n-2)\right\rfloor+1, & n\geq 20.
\end{cases}
\end{equation}
At $n=20$, the proposed value is $53$, whereas Euler's bound is $54$.

\begin{theorem}\label{thm:main}
The $n=20$ case of \eqref{eq:conjecture} holds; namely,
\[
  \exP(20,P_6^{\ind})=53.
\]
\end{theorem}

No claim is made here for $n\geq 21$.

We describe the lower-bound graph without a drawing. Begin with a path
$c_0c_1c_2c_3c_4$. Add adjacent vertices $a,b$, and join each of them to
every $c_i$. For $0\leq i\leq 3$, insert $u_i$ into the face
$ac_ic_{i+1}$ and $v_i$ into the face $bc_ic_{i+1}$; then insert $w_i$
into the face $u_ic_ic_{i+1}$. The resulting triangulation $T$ has
$19$ vertices and $51$ edges.

Place a new vertex $x$ in the face $abc_0$ and add the two edges $xa$ and
$xb$, but not $xc_0$. The resulting planar graph $W$ has $20$ vertices
and $53$ edges. It is the $n=20$ member of the construction in
\cite[Figure~5 and the preceding paragraph]{GyoriKarim}.

The triangulation $T$ is the one drawn in that figure, which the source
asserts, without proof, to contain no induced $P_6$;
Proposition~\ref{prop:census} below confirms the assertion. The source also
observes that repeatedly deleting a degree-three vertex from it yields an
induced-$P_6$-free triangulation for every smaller $n$, so the $n\leq 19$
branch of \eqref{eq:conjecture} is already settled there.

\section{Computational census}

Let $\cL_n$ be a set containing one representative from each isomorphism
class of induced-$P_6$-free triangulations on $n$ vertices.

\begin{lemma}\label{lem:m4}
No simple triangulation of minimum degree at least four on $n$ vertices,
$12\leq n\leq 20$, is induced-$P_6$-free.
\end{lemma}

The upper bound rests on Lemma~\ref{lem:m4} twice: it empties the
minimum-degree-four branch of the recursion below at every order from $12$ to
$19$, and it is the minimum-degree-four half of $\cL_{20}=\varnothing$. Its
content is a single sweep. For $12\leq n\leq 20$ the command
\texttt{plantri -m4 $n$} of \texttt{plantri} version~5.8 generates
$78{,}435{,}562$ isomorphism classes in total, and the induced-path test
described below rejected every one of them. The remaining orders of that
branch are small: for $6\leq n\leq 11$ the same command produces $55$ classes
altogether. The other branch is small as well: it tests $|\cL_{n-1}|(2n-6)$
facial insertions at order $n$, and $20{,}196$ candidates over
$5\leq n\leq 20$ in all.

\begin{proposition}\label{prop:census}
The graph $W$ is induced-$P_6$-free, and
\[
  \cL_{19}=\{T\},\qquad \cL_{20}=\varnothing.
\]
More specifically, no $20$-vertex triangulation of minimum degree at least
four is induced-$P_6$-free, and each of the $34$ facial insertions of $T$
contains an induced $P_6$.
\end{proposition}

\begin{proof}
For each $6\leq n\leq 20$, the minimum-degree-four branch was generated
with \texttt{plantri} version~5.8 \cite{Plantri,BrinkmannMcKay}, using
\[
  \texttt{plantri -m4 $n$}.
\]
The range starts at $6$ because there is no simple triangulation of minimum
degree at least four on fewer than six vertices: $K_4$ has all degrees three,
and the unique five-vertex triangulation has two degree-three vertices.
Since simple triangulations are $3$-connected, the generated plane-isomorphism
classes are also abstract isomorphism classes. The runs at orders $19$ and
$20$ produced, respectively,
\[
  11{,}284{,}042
  \qquad\text{and}\qquad
  64{,}719{,}885
\]
triangulations, in agreement with the published enumeration
\cite[Table~3]{BrinkmannMcKayFull}. At orders $19$ and $20$, none passed the
induced-$P_6$ test. Over the whole branch the test was passed only at orders
$6\leq n\leq 10$; for every $11\leq n\leq 20$ the branch contributed no
induced-$P_6$-free triangulation at all. This proves Lemma~\ref{lem:m4}.

The remaining triangulations were generated recursively. Every simple
triangulation has minimum degree at least three. If it has a degree-three
vertex, deleting that vertex leaves a smaller triangulation; conversely,
the original graph is recovered by inserting a vertex into a triangular
face. Since induced-$P_6$-freeness is hereditary under vertex deletion,
starting from $\cL_4=\{K_4\}$, adjoining the induced-$P_6$-free
minimum-degree-four outputs and the induced-$P_6$-free facial extensions of
members of $\cL_{n-1}$, followed by canonical-label isomorphism
rejection, yields exactly $\cL_n$. The two branches are disjoint, since a
triangulation of minimum degree at least four has no degree-three vertex and
so is not a facial extension of a smaller triangulation. Consequently
$|\cL_n|$ is the number of induced-$P_6$-free classes in the output of
\texttt{plantri -m4 $n$} plus the number of induced-$P_6$-free classes
obtained by facial extension from $\cL_{n-1}$, and by the previous paragraph
the first summand vanishes for $11\leq n\leq 20$. The computation gives
\[
  \bigl(|\cL_n|\bigr)_{n=4}^{20}
  =(1,1,2,5,13,36,77,126,169,190,154,101,47,20,4,1,0).
\]
The unique class at order $19$ is isomorphic to $T$.

For the induced-path test, starting from every oriented edge, a partial path
$z_1\ldots z_t$ was extended only to a vertex
\[
  y\in N(z_t)\setminus\{z_1,\ldots,z_t\}
  \quad\text{with}\quad
  N(y)\cap\{z_1,\ldots,z_{t-1}\}=\varnothing.
\]
Reaching six vertices is therefore equivalent to finding an induced
$P_6$. A separate six-subset implementation used the equivalent criterion
that the induced graph be connected, have five edges, and have degree
multiset $\{1,1,2,2,2,2\}$. The two implementations returned identical
results. The subset test also verifies directly that $W$ is
induced-$P_6$-free and supplies an induced path for each of the $34$ facial
insertions of $T$.
\end{proof}

\begin{proof}[Proof of Theorem~\ref{thm:main}]
Proposition~\ref{prop:census} gives the lower bound
$\exP(20,P_6^{\ind})\geq 53$ through the graph $W$.

Suppose that an induced-$P_6$-free planar graph $G$ on $20$ vertices has
$54$ edges. Equality in Euler's bound forces $G$ to be a triangulation.
If $\delta(G)\geq 4$, this contradicts Proposition~\ref{prop:census}.
Otherwise, $G$ has a degree-three vertex $v$. Then $G-v$ is an
induced-$P_6$-free triangulation on $19$ vertices, and hence $G-v\cong T$.
A simple triangulation is $3$-connected and has a unique spherical embedding
up to reflection, so $G$ is obtained by inserting $v$ into one of the
$34$ faces of $T$. This again contradicts
Proposition~\ref{prop:census}. Therefore $54$ edges are impossible, and
the upper bound is $53$.
\end{proof}

\section*{Reproducibility}

The census used \texttt{plantri} version~5.8 \cite{Plantri} and the two
implementations of the induced-$P_6$ test described in the proof of
Proposition~\ref{prop:census}. The source code, the \texttt{plantri} source
used and its SHA-256 hash, the compilation and execution commands, the logs,
and an induced $P_6$ for each of the $34$ facial insertions of $T$ are to be
deposited as ancillary files; none of them is needed for the argument above.

The two graphs the reader may wish to recheck are recorded here directly. In
\texttt{graph6} format, with the vertices in the order
$a$, $b$, $c_0,\ldots,c_4$, $u_0,\ldots,u_3$, $v_0,\ldots,v_3$,
$w_0,\ldots,w_3$, and $x$ last,
\begin{center}
\verb+R~vMLaWbCKM?U?R?Go@a?EG?KO?KO?+
\end{center}
encodes $T$, and
\begin{center}
\verb+S~vMLaWbCKM?U?R?Go@a?EG?KO?KOE???+
\end{center}
encodes $W$.

\begin{thebibliography}{9}

\bibitem{BrinkmannMcKay}
G.~Brinkmann and B.~D.~McKay,
\emph{Fast generation of planar graphs},
MATCH Commun. Math. Comput. Chem. \textbf{58} (2007), no.~2, 323--357.

\bibitem{BrinkmannMcKayFull}
G.~Brinkmann and B.~D.~McKay,
\emph{Fast generation of planar graphs}, expanded edition, which contains
tables absent from the journal version,
\url{https://users.cecs.anu.edu.au/~bdm/papers/plantri-full.pdf}.

\bibitem{Plantri}
G.~Brinkmann, H.~Van den Camp, and B.~D.~McKay,
\emph{plantri}, version 5.8, 2026,
\url{https://users.cecs.anu.edu.au/~bdm/plantri/}.

\bibitem{GyoriKarim}
E.~Gy\H{o}ri and H.~Hama Karim,
\emph{Induced planar Tur\'an numbers},
arXiv:2604.25829 [math.CO], 2026,
\url{https://doi.org/10.48550/arXiv.2604.25829}.

\end{thebibliography}

\end{document}